\documentclass[10pt,a4paper]{article}
\usepackage[utf8]{inputenc}
\usepackage{gdmath}

\usepackage[margin=2.4cm]{geometry}

\title{Some Proofs in Category Theory: The Yoneda Lemma}
\author{Georgy Dunaev \\ \texttt{georgedunaev@gmail.com}}
\date{}

\begin{document}
\sloppy
\maketitle
\section{Logic and Set Theory}
\begin{table}[H]
\centering
\setlength{\tabcolsep}{4.5pt}
\caption{Axiom schemata of predicate calculus with equality}
%\caption{Схемы аксиом логики предикатов с равенством}
\label{tab:axiompltab}
\begin{tabular}{|c|c|c|c|}
\hline
% назвывние & утверждение & условие\\ \hline
name & statement & condition & parameters(meta)\\ \hline
\hline
impI & $A \to (B \to A)$ & - & $A,B \in \textrm{Fm}$\\ %\hline
impE & $(A \to (B \to C))\to((A\to B)\to(A\to C))$ & -& $A,B,C \in \textrm{Fm}$\\ \hline
\hline
andE1 & $A \land B \to A$& -& $A,B \in \textrm{Fm}$\\ %\hline
andE2 & $A \land B \to B$& -& $A,B \in \textrm{Fm}$\\ %\hline
andI & $A\to(B\to (A \land B))$& -& $A,B \in \textrm{Fm}$\\ \hline
\hline
orI1 & $A \to  A \lor B$& -& $A,B \in \textrm{Fm}$\\ %\hline
orI2 & $B \to  A \lor B$& -& $A,B \in \textrm{Fm}$\\ %\hline
orE & $(A \to  C)\to ((B \to  C) \to  (A \lor B \to  C))$& -& $A,B,C \in \textrm{Fm}$\\ \hline
\hline
negI & $(A \to  B) \to  ((A \to  \neg B) \to  \neg A)$& -& $A,B \in \textrm{Fm}$\\ %\hline
negE & $A \to  (\neg A \to  B)$& -& $A,B \in \textrm{Fm}$\\ \hline
\hline
LEM & $A \lor \neg A$& -& $A \in \textrm{Fm}$\\ \hline
\hline
$\forall$E & $\forall xA\to A(t/x)$ & $\textrm{FFI}(t,x,A)$& $A \in \textrm{Fm}$\\ %\hline
$\exists$I & $A(t/x)\to \exists x A$ & $\textrm{FFI}(t,x,A)$& $A \in \textrm{Fm}$\\ \hline
\hline
=refl&$\forall w (w=w)$&-& $w \in \textrm{Var}$\\
=subs&$(x=y)\to ( A(x/y)\to A) $&$\textrm{FFI}(x,y,A)$& $A \in \textrm{Fm}$\\\hline
\end{tabular}
\end{table}
%Схема аксиом логики предикатов без равенства получается, если удалить последние две аксиомы.
%Здесь мы обозначили как $\textrm{FFI}(t,x,A)$ утверждение ``Терм $t$ свободен для подстановки вместо переменной $x$ в формуле $A$'', и как $A(t/x)$ -- результат корректной подстановки терма t вместо x.\\

\begin{table}[H]
\centering
\caption{Rules of inference of predicate calculus}
\label{tab:rulestab}
\begin{tabular}{|c|c|c|c|}
\hline
name & rule & condition & parameters(meta) \\ \hline
\hline
MP & \rule[-2.5ex]{0pt}{7ex} $\frac{A\to B\quad A}{B}$ & - & $A,B \in \textrm{Fm}$ \\ \hline
\hline
\rule[-2.5ex]{0pt}{7ex}
$\forall$I & $\frac{A\to B}{A\to \forall x B}$ & $x \notin \textrm{FV}(A)$ & $A,B \in \textrm{Fm},\ x \in \textrm{Var}$ \\ \hline
\hline
\rule[-2.5ex]{0pt}{7ex}
$\exists$E & $\frac{A\to B}{\exists x A\to B}$ & $x \notin \textrm{FV}(B)$ & $A,B \in \textrm{Fm},\ x \in \textrm{Var}$ \\ \hline
\end{tabular}
\end{table}

\subsection{Intersections}
Throughout this section, $\Set$ denotes the universal class of all sets (so $x \in \Set$ means that $x$ is a set).

\begin{definition}[Class intersection]
$$\interclass K \eqdef \{x \in \Set : \forall k\in K.\ x\in k\}.$$
\end{definition}

\begin{exercise}
$K\neq \emptyset \to \interclass K \in \Set$.
\end{exercise}
\myprf\ Since $K \neq \emptyset$, choose some $k_0 \in K$. Because $k_0$ is a set and every $x \in \interclass K$ satisfies $x \in k_0$, we have $\interclass K = \{x \in k_0 : \forall k\in K.\ x\in k\}$, which is a set by the Axiom Schema of Separation. \myqed

\begin{exercise}
$K=\emptyset \to \interclass K = \Set$.
\end{exercise}
\myprf\ If $K = \emptyset$, the condition $\forall k \in \emptyset.\ x \in k$ holds vacuously for every $x \in \Set$. Thus $\interclass \emptyset = \{x \in \Set : \top\} = \Set$ (a proper class). \myqed

\begin{definition}[Set intersection]
$$\inter K \eqdef \{x \in \Set : (\forall k\in K.\ x\in k)\land (\exists k.\ k\in K)\}.$$
\end{definition}

\begin{exercise}
$\inter K \in \Set$.
\end{exercise}
\myprf\ By the Law of Excluded Middle (LEM), either $K \neq \emptyset$ or $K = \emptyset$. If $K \neq \emptyset$, then $\exists k.\ k \in K$ holds, so $\inter K = \interclass K \in \Set$ by Exercise 1.1. If $K = \emptyset$, then $\exists k.\ k \in K$ is false, so $\inter \emptyset = \emptyset \in \Set$. \myqed

\subsection{Binders}
\begin{definition}[Uniqueness quantifier]
$$!x.\ P(x) \propdef \forall x_1\, \forall x_2.\ (P(x_1)\land P(x_2)\to x_1=x_2).$$
\end{definition}

\begin{definition}[``Exists and unique'' quantifier]
$$\exists !x.\ P(x) \propdef (\exists x.\ P(x))\land(!x.\ P(x)).$$
\end{definition}

\begin{definition}[Dual quantifier]
$$?x.\ P(x) \propdef \neg (!x.\ \neg P(x)).$$
\end{definition}

\begin{exercise}
$!x.\ P(x) \LRA \neg (?x.\ \neg P(x))$.
\end{exercise}
\myprf\ By definition of $?x$, $\neg (?x.\ \neg P(x)) \LRA \neg \neg (!x.\ \neg \neg P(x))$, which is equivalent in classical predicate logic (via LEM) to $!x.\ P(x)$. \myqed

\begin{definition}[Definite description binder]
$$\The x.\ P(x) \eqdef \inter \{x : P(x)\ \land\ !w.\ P(w)\}.$$
\end{definition}

\begin{exercise}
$(P(a)\ \land\ !w.\ P(w)) \to (\The x.\ P(x))=a$.
\end{exercise}
\myprf\ Assume $P(a)$ and $!w.\ P(w)$. For any $x$, $P(x) \land !w.\ P(w)$ holds iff $x = a$ (since $P(x) \land P(a) \to x=a$). Hence $\{x : P(x) \land !w.\ P(w)\} = \{a\}$, and $\inter \{a\} = \{z : z \in a\} = a$. \myqed

\begin{exercise}
$(\exists !w.\ P(w)) \to P(\The x.\ P(x))$.
\end{exercise}
\myprf\ Assume $\exists !w.\ P(w)$, i.e., $(\exists a.\ P(a)) \land (!w.\ P(w))$. By $\exists\textrm{E}$, pick $a$ such that $P(a)$. By Exercise 1.5, $(\The x.\ P(x)) = a$. By equality substitution ($=\textrm{subs}$), $P(\The x.\ P(x))$ holds. \myqed

\subsection{Values of functions}
For a relation $g$, its value at $x$ is defined via definite description:
$$g(x) \eqdef (g\apo x) \eqdef (\The y.\ \langle x, y\rangle \in g) = \inter \{y : \langle x,y\rangle \in g \land !w.\ \langle x,w\rangle \in g\}.$$

\begin{exercise}
$x\notin \dom(g)\to (g\apo x)=\emptyset$.
\end{exercise}
\myprf\ If $x \notin \dom(g)$, there is no $y$ with $\langle x,y\rangle \in g$. Hence $\{y : \langle x,y\rangle \in g \land !w.\ \langle x,w\rangle \in g\} = \emptyset$, and $\inter \emptyset = \emptyset$. \myqed

\subsection{Dependent sum and product}
Let $R$ be a family of sets indexed by a set (or class) $W = \dom(R)$, so that $R(x) \in \Set$ for each $x \in W$.
\begin{definition}[Dependent sum]
$$\Sigma R \eqdef \{\langle x, y\rangle : x \in \dom(R) \land y\in R(x)\}.$$
\end{definition}
\begin{definition}[Dependent product]
$$\Pi R \eqdef \{ f\in \Pow(\Sigma R) : \Rel(f) \land \Fun(f) \land \dom(R)\subseteq \dom(f) \land \forall x\in\dom(R).\ f(x)\in R(x)\}.$$
When $R(x)$ is written as $B(x)$ for $x \in A$, we also write $\Pi(x \in A).\ B(x)$ for $\Pi R$.
\end{definition}

\subsection{Extensionality}
Recall that $\Rel(g) \propdef \forall p \in g.\ \exists a\, \exists b.\ p = \langle a,b\rangle$, and $\Fun(g) \propdef \forall x.\ !y.\ \langle x,y\rangle \in g$.

\begin{theorem}[Inclusion of functions]
Assume $\Rel(g_1)$, $\Fun(g_1)$, $\Fun(g_2)$, $\dom(g_1)\subseteq\dom(g_2)\subseteq X$, and $\forall x\in X.\ g_1(x)=g_2(x)$. Then $g_1\subseteq g_2$.
\end{theorem}
\myprf
\begin{enumerate}
\item Let $p \in g_1$. Since $\Rel(g_1)$, there exist $a, b$ such that $p = \langle a, b\rangle \in g_1$.
\item Since $\langle a, b\rangle \in g_1$, we have $a \in \dom(g_1)$. Because $\Fun(g_1)$ holds, $!w.\ \langle a,w\rangle \in g_1$, so by Exercise 1.5, $g_1(a) = (\The y.\ \langle a,y\rangle \in g_1) = b$.
\item Since $a \in \dom(g_1) \subseteq \dom(g_2) \subseteq X$, the hypothesis $\forall x \in X.\ g_1(x) = g_2(x)$ yields $g_2(a) = g_1(a) = b$.
\item Moreover, $a \in \dom(g_2)$ and $\Fun(g_2)$ imply $\exists !y.\ \langle a,y\rangle \in g_2$. By Exercise 1.6, $\langle a, g_2(a)\rangle \in g_2$.
\item Substituting $g_2(a) = b$, we obtain $p = \langle a, b\rangle \in g_2$. Therefore $g_1 \subseteq g_2$.
\end{enumerate}
\myqed

\begin{theorem}[Functional extensionality]
Assume $\Rel(g_1)$, $\Rel(g_2)$, $\Fun(g_1)$, $\Fun(g_2)$, $\dom(g_1) = \dom(g_2) = X$, and $\forall x\in X.\ g_1(x)=g_2(x)$. Then $g_1 = g_2$.
\end{theorem}
\myprf\ Applying Theorem 1.1 to $(g_1, g_2)$ gives $g_1 \subseteq g_2$, and applying Theorem 1.1 symmetrically to $(g_2, g_1)$ gives $g_2 \subseteq g_1$. By set extensionality, $g_1 = g_2$. \myqed
\subsection{Transitive Sets and Ordinals}
\begin{definition}[Transitive set and successor]
$$\Tr(A) \propdef \forall x\in A.\ x\subseteq A, \qquad \mysucc(A) \eqdef A \cup \{A\}.$$
\end{definition}

\begin{definition}[Transitive ordinal predicate]
$$\TOrd(A) \propdef \Tr(A)\land \forall a\in A.\ \Tr(a).$$
\end{definition}

\begin{theorem}
$\exists A.\ \Tr(A) \land \neg (\forall a\in A.\ \Tr(a))$.
\end{theorem}
\myprf\ Let $0 \eqdef \emptyset$ and $1 \eqdef \{0\} = \{\emptyset\}$. (Note that $\{1\} = \{\{\emptyset\}\}$ is not transitive because $0 \in 1 \in \{1\}$ but $0 \notin \{1\}$.) Define
$$A \coloneqq \{0, 1, \{1\}\} = \{\emptyset, \{\emptyset\}, \{\{\emptyset\}\}\}.$$
\begin{enumerate}
\item $\Tr(A)$ holds: the elements of $A$ are $0 = \emptyset \subseteq A$, $1 = \{0\} \subseteq A$, and $\{1\} \subseteq A$. Thus $\forall x \in A.\ x \subseteq A$.
\item $\neg (\forall a \in A.\ \Tr(a))$ holds: take $a = \{1\} \in A$. Since $1 \in \{1\}$ and $0 \in 1$, transitivity of $\{1\}$ would require $0 \in \{1\}$, which is false ($0 \neq 1$). Hence $\neg \Tr(\{1\})$.
\end{enumerate}
\myqed

\begin{theorem}
$\Tr(A) \longrightarrow \Tr(\mysucc(A))$.
\end{theorem}
\myprf\ Assume $\Tr(A)$, so $\forall x\in A.\ x\subseteq A$. Let $x \in \mysucc(A) = A\cup\{A\}$. Then $x\in A \lor x=A$:
\begin{enumerate}
\item Case $x\in A$: by $\Tr(A)$, $x\subseteq A \subseteq A\cup\{A\}$.
\item Case $x = A$: $x = A \subseteq A\cup\{A\}$.
\end{enumerate}
Hence $\Tr(A\cup\{A\})$. \myqed

\begin{theorem}
$\TOrd(A) \longrightarrow \TOrd(\mysucc(A))$.
\end{theorem}
\myprf\ From $H: \TOrd(A)$ we have $H_1: \Tr(A)$ and $H_2: \forall a\in A.\ \Tr(a)$.
\begin{enumerate}
\item Aim 1: $\Tr(\mysucc(A))$, which follows immediately from Theorem 1.4 and $H_1$.
\item Aim 2: $\forall x\in A\cup\{A\}.\ \Tr(x)$. For any $x \in A \cup \{A\}$:
  \begin{itemize}
  \item If $x\in A$, then $\Tr(x)$ holds by $H_2$.
  \item If $x = A$, then $\Tr(x)$ holds by $H_1$.
  \end{itemize}
\end{enumerate}
Thus $\TOrd(\mysucc(A))$. \myqed
\section{Category Theory}
\noindent\textbf{Notes and notation:}
\begin{itemize}
\item $\Ob_C$ and $C$ are used interchangeably for the class of objects of a category $C$ (e.g., $\Ob_\Set$ and $\Set$).
\item Composition of morphisms $g \in \Hom_C(X,Y)$ and $f \in \Hom_C(Y,Z)$ is written in applicative order:
$$\circ_C^{X,Y,Z} : \Hom_C(Y,Z) \times \Hom_C(X,Y) \to \Hom_C(X,Z), \qquad (f, g) \mapsto f \circ_C^{X,Y,Z} g.$$
Composition is associative: $h \circ_C (g \circ_C f) = (h \circ_C g) \circ_C f$, and has identities $1_X = \id_X \in \Hom_C(X,X)$ satisfying $f \circ_C 1_X = f$ and $1_Y \circ_C f = f$ for all $f \in \Hom_C(X,Y)$.
\item \textbf{Duality:} The opposite category $C^\op$ has $\Ob_{C^\op} = \Ob_C$ and $\Hom_{C^\op}(A,B) = \Hom_C(B,A)$ for all $A,B \in \Ob_C$. For $f \in \Hom_{C^\op}(Y,Z) = \Hom_C(Z,Y)$ and $g \in \Hom_{C^\op}(X,Y) = \Hom_C(Y,X)$, composition in $C^\op$ is given by
$$f \circ_{C^\op}^{X,Y,Z} g \eqdef g \circ_C^{Z,Y,X} f \in \Hom_C(Z,X) = \Hom_{C^\op}(X,Z).$$
\end{itemize}

\subsection{Yoneda Embedding and the Yoneda Lemma}
Let $C$ be a locally small category (meaning $\Hom_C(A,B) \in \Set$ for all $A,B \in \Ob_C$), and let $\Set$ be the category of sets.

\begin{definition}[Covariant functor]
Let $C$ and $D$ be categories. A pair $(F, \FMor)$ is a (covariant) functor from $C$ to $D$ iff:
\begin{enumerate}
\item $F : \Ob_C \to \Ob_D$,
\item $\FMor : \Pi(A, B \in \Ob_C).\ \Hom_C(A,B) \to \Hom_D(F(A), F(B))$,
\item $\forall A\in \Ob_C.\ \FMor^{A,A}(1_A)=1_{F(A)}$,
\item $\forall X,Y,Z \in \Ob_C.\ \forall f\in\Hom_C(Y,Z).\ \forall g\in\Hom_C(X,Y).$
$$\FMor^{X,Z}(f\circ_C^{X,Y,Z} g)=\FMor^{Y,Z}(f)\circ_D^{F(X),F(Y),F(Z)} \FMor^{X,Y}(g).$$
\end{enumerate}
A \emph{contravariant functor} from $C$ to $D$ is a covariant functor from $C^\op$ to $D$.
\end{definition}

\begin{definition}[Functor class]
Let $C, D$ be categories. We denote by $[C,D]$ the class (or category) of functors from $C$ to $D$.
\end{definition}

\begin{definition}[Morphism map of the contravariant Hom-functor]
Fix an object $x \in \Ob_C$. On objects, $\Hom_C(-,x) : \Ob_{C^\op} \to \Set$ sends $A \mapsto \Hom_C(A,x)$.
For $A,B \in \Ob_{C^\op}$ and $f \in \Hom_{C^\op}(A,B) = \Hom_C(B,A)$, we define the morphism map (pre-composition with $f$):
$$\HomMor^{A,B}_C(f,x) : \Hom_C(A,x) \to \Hom_C(B,x),$$
$$\HomMor^{A,B}_C(f,x) \eqdef \left(g \in \Hom_C(A,x) \mapsto (g \circ_C^{B,A,x} f) \in \Hom_C(B,x)\right).$$
(An alternative notation for $\HomMor^{A,B}_C(f,x)$ is $\widetilde{\Hom_C(-,x)}^{A,B}(f)$.)
\end{definition}

\begin{lemma}[Contravariant Hom-functor]
For every $x \in \Ob_C$, $\left(\Hom_C(-, x),\, \HomMor_C(-,x)\right) \in [C^\op,\Set]$.
\end{lemma}
\myprf\ We verify the four conditions of Definition 2.1 with $C \coloneqq C^\op$, $D \coloneqq \Set$, $F \coloneqq \Hom_C(-,x)$, and $\FMor \coloneqq \HomMor_C(-,x)$:
\begin{enumerate}
\item \textbf{Object map:} $\Hom_C(-, x) : \Ob_{C^\op} \to \Set$ since $C$ is locally small ($\Hom_C(A,x) \in \Set$ for every $A \in \Ob_C$).
\item \textbf{Morphism map:} Fix $A,B \in \Ob_{C^\op}$ and $f \in \Hom_{C^\op}(A,B)$. By definition of $C^\op$, $f \in \Hom_C(B,A)$. Thus for every $g \in \Hom_C(A,x)$, the composite $g \circ_C^{B,A,x} f$ is a well-defined element of $\Hom_C(B,x)$. Hence
$$\HomMor_C^{A,B}(f, x) = \left(g \in \Hom_C(A,x) \mapsto g \circ_C f \in \Hom_C(B,x)\right) \in \Hom_\Set(\Hom_C(A,x), \Hom_C(B,x)).$$
\item \textbf{Preservation of identity:} Let $A \in \Ob_{C^\op}$ and $1_A \in \Hom_{C^\op}(A,A) = \Hom_C(A,A)$. Using the right identity law $g \circ_C 1_A = g$ for every $g \in \Hom_C(A,x)$, we obtain
\begin{align*}
\HomMor_C^{A,A}(1_A, x) &= \left(g \in \Hom_C(A,x) \mapsto g \circ_C 1_A\right) \\
&= \left(g \in \Hom_C(A,x) \mapsto g\right) = 1_{\Hom_C(A,x)}.
\end{align*}
\item \textbf{Preservation of composition:} Fix $X,Y,Z \in \Ob_{C^\op}$, $f \in \Hom_{C^\op}(Y,Z) = \Hom_C(Z,Y)$, and $g \in \Hom_{C^\op}(X,Y) = \Hom_C(Y,X)$. We must show:
$$\HomMor_C^{X,Z}(f \circ_{C^\op}^{X,Y,Z} g, x) = \HomMor_C^{Y,Z}(f,x) \circ_\Set^{\Hom_C(X,x),\Hom_C(Y,x),\Hom_C(Z,x)} \HomMor_C^{X,Y}(g,x).$$
Since $f \circ_{C^\op}^{X,Y,Z} g = g \circ_C^{Z,Y,X} f \in \Hom_C(Z,X)$, we compare the two maps from $\Hom_C(X,x)$ to $\Hom_C(Z,x)$:
\begin{enumerate}
\item[(a)] $\HomMor_C^{X,Z}(g \circ_C f, x) = \left(h \in \Hom_C(X,x) \mapsto h \circ_C (g \circ_C f) \in \Hom_C(Z,x)\right)$,
\item[(b)] $\HomMor_C^{X,Y}(g, x) = \left(h \in \Hom_C(X,x) \mapsto (h \circ_C g) \in \Hom_C(Y,x)\right)$,
\item[(c)] $\HomMor_C^{Y,Z}(f, x) = \left(v \in \Hom_C(Y,x) \mapsto (v \circ_C f) \in \Hom_C(Z,x)\right)$.
\end{enumerate}
Composing (b) and (c) in $\Set$ along $\Hom_C(X,x) \xrightarrow{\HomMor_C^{X,Y}(g, x)} \Hom_C(Y,x) \xrightarrow{\HomMor_C^{Y,Z}(f, x)} \Hom_C(Z,x)$ acts on $h \in \Hom_C(X,x)$ as:
$$h \xmapsto{\HomMor_C^{X,Y}(g, x)} (h \circ_C g) \xmapsto{\HomMor_C^{Y,Z}(f, x)} (h \circ_C g) \circ_C f.$$
By associativity of composition in $C$, $(h \circ_C g) \circ_C f = h \circ_C (g \circ_C f)$ for all $h \in \Hom_C(X,x)$. By functional extensionality (Theorem 1.2), $\HomMor_C^{X,Z}(f \circ_{C^\op} g, x) = \HomMor_C^{Y,Z}(f,x) \circ_\Set \HomMor_C^{X,Y}(g,x)$.
\end{enumerate}
\myqed
\begin{definition}[Natural transformation and functor category]
Let $C, D$ be categories and let $(F, \FMor), (G, \GMor) \in [C, D]$ be covariant functors. A \emph{natural transformation} $\eta$ from $F$ to $G$, written $\eta \in \Nat(F, G)$, is a dependent family of morphisms
$$\eta : \Pi(a \in \Ob_C).\ \Hom_D(F(a), G(a))$$
satisfying the \emph{naturality condition}:
$$\forall a, b \in \Ob_C.\ \forall m \in \Hom_C(a,b).\quad \eta_b \circ_D \FMor^{a,b}(m) = \GMor^{a,b}(m) \circ_D \eta_a.$$
The functor category $[C, D]$ has objects $\Ob_{[C,D]} = [C,D]$ and morphism sets $\Hom_{[C,D]}(F,G) = \Nat(F,G)$, with identity $\left(1_F\right)_a = 1_{F(a)}$ and objectwise composition $(\theta \circ_{[C,D]} \eta)_a = \theta_a \circ_D \eta_a$.
\end{definition}

\begin{definition}[Representable presheaf $h_x$]
For each object $x \in \Ob_C$, the representable presheaf $h_x \in \Ob_{[C^\op,\Set]}$ is the contravariant Hom-functor
$$h_x \eqdef \Hom_C(-, x) : C^\op \to \Set.$$
\end{definition}

\begin{definition}[Induced family $\eta^f$ by post-composition]
Let $x, y \in \Ob_C$ and $f \in \Hom_C(x,y)$. Define the family of maps
$$\eta^f : \Pi(w\in \Ob_{C^\op}).\ \Hom_C(w,x)\to \Hom_C(w,y),$$
$$\eta^f_w \eqdef \left(\lambda g\in\Hom_C(w,x).\ f\circ_C^{w,x,y} g\right).$$
\end{definition}

\begin{lemma}[Naturality of $\eta^f$]
Let $x,y \in \Ob_C$ and $f \in \Hom_C(x,y)$. Then $\eta^f$ is a natural transformation from $h_x = \Hom_C(-,x)$ to $h_y = \Hom_C(-,y)$, i.e., $\eta^f \in \Nat(h_x, h_y)$:
$$\forall a,b\in \Ob_{C^\op}.\ \forall m\in \Hom_{C^\op}(a,b).\quad \eta^f_b \circ_\Set \HomMor_C^{a,b}(m,x) = \HomMor_C^{a,b}(m,y)\circ_\Set\eta^f_a.$$
\end{lemma}
\myprf\ Fix $a, b \in \Ob_{C^\op}$ and $m \in \Hom_{C^\op}(a,b) = \Hom_C(b,a)$. We compare the two composite functions in $\Hom_\Set(\Hom_C(a,x), \Hom_C(b,y))$:
\begin{enumerate}
\item By Definition 2.3, $\HomMor_C^{a,b}(m,x) : \Hom_C(a,x) \to \Hom_C(b,x)$ is $\lambda g \in \Hom_C(a,x).\ g \circ_C^{b,a,x} m$, and $\eta^f_b : \Hom_C(b,x) \to \Hom_C(b,y)$ is $\lambda g_1 \in \Hom_C(b,x).\ f \circ_C^{b,x,y} g_1$. Thus the left-hand side is:
$$\textrm{LHS} = \eta^f_b \circ_\Set \HomMor_C^{a,b}(m,x) = \lambda g \in \Hom_C(a,x).\ f \circ_C^{b,x,y} \left(g \circ_C^{b,a,x} m\right).$$
\item Similarly, $\eta^f_a : \Hom_C(a,x) \to \Hom_C(a,y)$ is $\lambda g \in \Hom_C(a,x).\ f \circ_C^{a,x,y} g$, and $\HomMor_C^{a,b}(m,y) : \Hom_C(a,y) \to \Hom_C(b,y)$ is $\lambda g_3 \in \Hom_C(a,y).\ g_3 \circ_C^{b,a,y} m$. Thus the right-hand side is:
$$\textrm{RHS} = \HomMor_C^{a,b}(m,y) \circ_\Set \eta^f_a = \lambda g \in \Hom_C(a,x).\ \left(f \circ_C^{a,x,y} g\right) \circ_C^{b,a,y} m.$$
\item By associativity of composition $\circ_C$ in $C$, $f \circ_C (g \circ_C m) = (f \circ_C g) \circ_C m$ for all $g \in \Hom_C(a,x)$. By functional extensionality (Theorem 1.2), $\textrm{LHS} = \textrm{RHS}$.
\end{enumerate}
\myqed
\begin{lemma}[Functoriality of the Yoneda embedding]
Define $\left(h_-, \widetilde{h}_-\right)$ from $C$ to $[C^\op, \Set]$ by $h_x \eqdef \Hom_C(-,x)$ on objects $x \in \Ob_C$ and $\widetilde{h}^{x,y}(f) \eqdef \eta^f$ on morphisms $f \in \Hom_C(x,y)$. Then $\left(h_-, \widetilde{h}_-\right) : C \to [C^\op, \Set]$ is a covariant functor.
\end{lemma}
\myprf\ We check the four conditions of Definition 2.1 with $C \coloneqq C$, $D \coloneqq [C^\op,\Set]$, $F \coloneqq h_-$, and $\FMor \coloneqq \widetilde{h}_-$:
\begin{enumerate}
\item \textbf{Object map:} For every $x \in \Ob_C$, $h_x = \Hom_C(-,x) \in \Ob_{[C^\op,\Set]}$ by Lemma 2.1.
\item \textbf{Morphism map:} For every $A,B \in \Ob_C$ and $f \in \Hom_C(A,B)$, $\widetilde{h}^{A,B}(f) = \eta^f \in \Nat(h_A, h_B) = \Hom_{[C^\op,\Set]}(h_A, h_B)$ by Lemma 2.2.
\item \textbf{Preservation of identity:} Let $A \in \Ob_C$. For any $w \in \Ob_{C^\op}$ and $g \in h_A(w) = \Hom_C(w,A)$,
$$\left(\widetilde{h}^{A,A}(1_A)\right)_w(g) = \eta^{1_A}_w(g) = 1_A \circ_C^{w,A,A} g = g = \left(1_{h_A}\right)_w(g)$$
by the left identity law in $C$. By functional extensionality, $\widetilde{h}^{A,A}(1_A) = 1_{h_A}$.
\item \textbf{Preservation of composition:} Let $X,Y,Z \in \Ob_C$, $f \in \Hom_C(Y,Z)$, and $g \in \Hom_C(X,Y)$. For any $w \in \Ob_{C^\op}$ and $k \in h_X(w) = \Hom_C(w,X)$,
\begin{align*}
\left(\widetilde{h}^{X,Z}(f \circ_C g)\right)_w(k) &= \eta^{f \circ_C g}_w(k) = (f \circ_C g) \circ_C k \\
&= f \circ_C (g \circ_C k) = \eta^f_w\left(\eta^g_w(k)\right) = \left(\widetilde{h}^{Y,Z}(f) \circ_{[C^\op,\Set]} \widetilde{h}^{X,Y}(g)\right)_w(k)
\end{align*}
by associativity of $\circ_C$. By functional extensionality, $\widetilde{h}^{X,Z}(f \circ_C g) = \widetilde{h}^{Y,Z}(f) \circ_{[C^\op,\Set]} \widetilde{h}^{X,Y}(g)$.
\end{enumerate}
\text{(Fullness, faithfulness, and injectivity of $h_-$ will be proved right after the Yoneda Lemma.)}
\myqed
\begin{lemma}{by N.~Yoneda. (Contravariant version)}\\
$G : C^\op \to \Set$\\
$h_A = \Hom_C(-,A)$\\
$h_A : C^\op \to \Set$\\
$\Nat(h_A,G)\sim G(A)$\\
(there is a bijection between value of functor G and natural transformations from $h_A$ to $G$)
\end{lemma}
\myprf\\
Let $u\in G(A)$. Let's define a natural transformation $\Phi$ from $h_A$ to $G$.\\
$\Phi:\Pi(x\in C^\op).h_A(x)\to G(x)$ (because it's a Hom in Set)\\
1) Let $\Phi^u_A(\id_A)\eqdef u$. \\
Now we need to define $\Phi_x(f)$ for arbitrary $x\in C$ and\\
$f\in h_A(x).$ $=\Hom_C(x,A).$ $=\Hom_{C^\op}(A,x). \Longrightarrow$\\
$\Longrightarrow$ therefore $\GMor ^{A,x}(f):G(A)\to G(x)$ is defined.\\
So the definition will be $\Phi^u_x(f)\eqdef (\GMor^{A,x}(f))(u)$.\\
Now we need to show the correctness of the definition.\\
$\Phi^u_A(\id_A)\eqdef (\GMor^{A,A}(f))(\id_A) = (\id_{G(A)})(u) = u$. The correctness is proved.
2) The next aim is proving commutativity of natural transformation with functor applied to morphisms.\\
$\forall x,y\in C^\op.\forall f\in \Hom_{C^\op}(x,y).(G(f)\circ_\Set \Phi_x) = (\Phi_y \circ_\Set (\Hom_C(f,A)))$
Fix  $x,y\in C^\op$ and $f\in \Hom_{C^\op}(x,y)$.\\
Aim: $(G(f)\circ_\Set \Phi_x) = (\Phi_y \circ_\Set (\Hom_C(f,A)))$.
Fix $g\in \Hom_C(x,A)$. Recall that $\Hom_C(f,A)=(-\circ^{y,x,a})$.\\
Aim: $(G(f)\circ_\Set \Phi_x)(g) = (\Phi_y \circ_\Set (\Hom_C(f,A)))(g)$.\\
Equivalently, $(G(f)\circ_\Set \Phi_x)(g) = \Phi_y (g\circ_C^{y,x,A} f)$.\\
LHS $\coloneqq (G(f)\circ_\Set \Phi_x)(g) =($def of $\Phi)= \GMor^{x,y}(f)\left(\left(\GMor^{A,x}\left(g\right)\right)\left(u\right)\right) =$\\
$ =\GMor^{A,y}\left(g\circ_C f\right)(u)$\\
Then since $g\circ_C f = f\circ_{C^\op} g$ and $\GMor(f\circ_{C^\op} g)=  \GMor(f)\circ_\Set \GMor(g)$ we have:\\
RHS $\coloneqq \Phi^u_y(g\circ f)=($def of $\Phi)=\GMor^{A,y}\left(g\circ_C f\right)(u)$\\
This LHS=RHS holds therefore we obtained a natural transformation $\Phi\in\Nat(h_A,G)$ using the $u\in G(A)$.\\
We already know that the relation between $G(A)$ and $\Phi\in\Nat(h_A,G)$ is functional and total on $G(A)$.
The next step is to prove its injectivity and surjectivity on $\Phi\in\Nat(h_A,G)$.\\
Injectivity: $u_1,u_2\in G(A)$, $u_1\neq u_2$.\\
$\Phi_{1,A}(\id(A))= u_1\neq u_2 = \Phi_{2,A}(\id(A))\quad \Longrightarrow \quad \Phi_1 \neq \Phi_2$. It's proved. \\
Surjectivity on $\Nat(h_A,G)$.:\\
Let $\Psi\in\Nat(h_A,G)$. Let's show that $\exists u\in G(A).\Phi^u=\Psi$.\\
$u\coloneqq \Psi_A(\id_A)$\\
$\Psi:\Pi(x\in C). h_A(x)\to G(x)$\\
$\Phi^u_x(f) \eqdef (\GMor^{A,x}(f))(u)$.\\
Fix $x\in C$ and $f\in\Hom_C(x,A)$. $\Psi_x(f)=\ ?$.
In other words, there is a $\Psi\in\Nat(h_A,G)$, and it is known that $\Psi_A(\id_A) = u$. Find $\Psi_x(f)$.\\
Firstly, $\id_A \in \Hom_C(A,A)$. $\Longrightarrow (\id_A \circ_C^{x,A,A} f) \in \Hom_C(x,A)$. $\Longrightarrow \Psi_x(\id_A \circ_C^{x,A,A} f)\in G(A)$. $\Longrightarrow \Psi_x(f)\in G(A)$.\\
Secondly, $\id_A \in \Hom_C(A,A)$. $\Longrightarrow \Psi_A(\id_A) \in G(A)$. Since $\Psi_A(\id_A)=u$ we have $G(f)(u)\in G(x)$.\\
Since $\Psi$ is a natural transformation we obtained that $\Psi_x(f)=G(f)(u),$ which is equal to $\Phi^u_x(f)$.\\
This equality holds for any argument $f\in \Hom_C(x,A)$ then $\forall x\in C. \Psi_x=\Phi^u_x$ by functional extensionality. Therefore the relation is surjective on $\Nat(h_A,G)$.
\myqed\\
\end{document}